\section{Conclusion}
\label{sec:conclusion}

% CONSTRAINTS THIS SECTION IS HELD TO, since a conclusion is where they slip:
%   - No number appears here that does not appear earlier. The figures used are
%     15 of 61, the ten most-annotated classes, five detectors, and the 45/46
%     evaluated-class counts, all from Sections 3 and 5.
%   - The 45 and 46 in the opening sentence MUST keep their "computed on the
%     published split" scoping. Unscoped they read as properties of the dataset
%     rather than of the partition, which is the error the whole paper is about.
%     The same scoping is carried at every other site: 3.1, 5.4, 7.1 and 7.2.
%   - No citation. Every sentence here restates something this paper
%     established, so a citation would either be redundant or would smuggle in a
%     claim the body has not made.
%   - No verbatim overlap with Section 7. Checked mechanically, the same way
%     Section 2 was: no sentence over 60 characters is shared between them.
%
% THE FIRST PARAGRAPH CARRIES THE DEVIATION, CHANGED 2026-08-15. It read "under
% a single configuration" and then "the benchmark holds every setting but the
% architecture fixed" -- two flattenings of 5.3 in consecutive sentences, the
% second flatly contradicted by it. RT-DETR-l needs a learning rate a
% hundredfold lower and does not train at all without it. Both now carry the
% exception, and "otherwise" is load-bearing in the second. Matched to 5.2's own
% wording and to the same edit in the abstract and Section 1's contribution 3.
%
% THE LAST PARAGRAPH DEPENDS ON SECTION 10. "released as a manifest, with the
% scripts and run records" has somewhere to point only because
% sections/10-availability.tex names a repository and an archived identifier.
% That DOI is a PLACEHOLDER as of 2026-08-15; if it is never minted, this
% paragraph and the abstract's closing sentence both need softening.
%
% The paragraph most at risk is the fourth. Section 7.2 lists three checks that
% would have CAUGHT the defect; this lists four things a study can REPORT. The
% overlap is real and deliberate -- they are the same underlying quantities --
% so the wording is kept deliberately distinct, and the framing is prospective
% rather than retrospective.

A detection result reported on ElectroCom61 measures less than its name suggests:
computed on the published split it averages over 45 classes on validation and 46
on test rather than 61, reported as a bare mAP@50 it rests on the metric whose
ranking moves when the split changes, and reported without a timing measurement
it cannot say whether the model is fast enough for the conveyor the dataset was
assembled for.

This paper audited ElectroCom61 from its released files alone, constructed and
released a corrected split, and trained five detectors on both the published and
the corrected partition under one configuration, with a single documented
deviation: the transformer detector will not train at the learning rate the four
convolutional ones are trained at. The audit needed no model and no pixels; the
benchmark otherwise holds every setting but the architecture fixed, so that
differences between the two splits are attributable to the split.

The audit found that the published partition cannot evaluate 15 of the 61
classes at all, because they hold no instances in either the validation or the
test set. The cause is how capture sessions were allocated rather than how rare
those classes are. Rarity would leave the least frequent classes unevaluable and
no others; instead the fifteen span almost the entire frequency range, and eight
of the ten most heavily annotated classes are among them.

The benchmark found that correcting the split raises every model's score, on
both accuracy metrics and both evaluation partitions, without exception. The
defect had been understating measured performance rather than overstating it,
because what it withheld from evaluation was easier than what it left in. That
is the reverse of what a defective split is normally expected to do, and it is
the reason such a defect is easy to leave in place: a result that is too low
attracts no scrutiny.

Four things could be reported at no cost by any study on any dataset, and each
would have exposed something the record for this one does not show. How many
classes the metric was actually computed over, which the framework prints
without being asked. How each capture session was divided, not only the totals
across the dataset. Which version of the data, which partition and which seed
produced a figure. And a timing measurement, without which an accuracy ranking
cannot say whether the better model is fast enough to use.

The corrected split is released as a manifest, with the scripts and run records
that produced and verified it, so that every figure reported here can be
recomputed rather than taken on trust. What this study cannot establish from one
dataset is how often the defect it describes occurs elsewhere; establishing that
would take the checks above, which cost minutes each, run against other
benchmarks whose collection was organised in sessions.
